%% T2: same weights, only the corrector changes. Filled from outputs/paper/builtin_opv2v.json
\begin{table}[t]
  \centering\small
  \caption{\textbf{T2: replacing built-in correction, same weights.} OPV2V, AP$_{0.7}$, 3 paired seeds.}
  \begin{tabular}{@{}llcccc@{}}
    \toprule
    Network & Corrector & 0 & 0.4 & 1 & 2 \\
    \midrule
    \multirow{4}{*}{CoAlign} & none & \pending & \pending & \pending & \pending \\
     & its pose graph & \pending & \pending & \pending & \pending \\
     & FreeAlign & \pending & \pending & \pending & \pending \\
     & \textbf{AlignFormer} & \pending & \pending & \pending & \pending \\
    \midrule
    \multirow{4}{*}{RoCo} & none & \pending & \pending & \pending & \pending \\
     & its matching & \pending & \pending & \pending & \pending \\
     & FreeAlign & \pending & \pending & \pending & \pending \\
     & \textbf{AlignFormer} & \pending & \pending & \pending & \pending \\
    \bottomrule
  \end{tabular}
  \why{\textbf{Why this table:} the weights are identical in all four rows of a network, so any AP difference comes from the corrector alone. It is the cleanest comparison in the paper. \textbf{How to read it.} ``None'' against the built-in row: what the network's own correction is worth. AlignFormer against the built-in row: whether a general module can replace it. \textbf{What we conclude.} If AlignFormer is at or above the built-in row: a network-specific correction stage is not needed. If it is below: built-in correction still has an edge, and we state the trade-off: slightly less AP, but one module for any host.}
\end{table}
